\chapter{Synthesis: evidence for the open decisions}

\noindent The Decision Log (\texttt{00\_Admin/Decision-Log.md}) lists twelve decisions that the methodology chapter may not state until each has a cited reason. This chapter maps each decision to the chapters and papers of this digest. It lists the evidence and the competing options found in the literature; it does not make the decision.

\begin{longtable}{@{}P{0.9cm}P{3.2cm}P{4.5cm}P{4.9cm}@{}}
\toprule
\textbf{ID} & \textbf{Question} & \textbf{Evidence in the digest} & \textbf{Open point}\\
\midrule
\endhead
D1 & Sample size and stopping rule & Chapter~A: 20 apps is the mode in app reviews (\pa{benjumea2020}); designs range from 20 (\pa{malki2024}) to 192 (\pa{schultchen2021}); census designs in tens of thousands; scoping reviews allow criteria refined after seeing the volume if logged (\pa{arksey2005}) & ``First 50 eligible'' needs the order of ranking, screening and stopping fixed in advance and the shortfall rule (\pa{page2021})\\
D2 & Store search, ranking, rating or download floor & Chapter~A: anonymous account and dated search (\pa{huckvale2019}, \pa{lin2025}); data-derived rating floor (\pa{alfawzan2022}); recent-review ranking (\pa{malki2024}); store rank varies by account and device (\pa{messner2020}, \pa{robillard2019}) & API order versus visible store rank is unverified; floor and ranking measure rating, not popularity (\pa{hutton2018})\\
D3 & Tracker scope and keywords & Chapter~S and A: function over category (\pa{robillard2019}, \pa{dehling2015}); sleep boundary (\pa{simon2023}); keywords from user data (\pa{hutton2018}) & Mental-health and specialised trackers need inclusion criteria in functional terms\\
D4 & Payment, trial, account eligibility & Chapter~A and D: exclusion with reasons (\pa{huckvale2015}); DiGA rules on in-app purchases (\pa{bfarm2023}) & Which paid or trial apps can be inspected under study accounts\\
D5 & Apps without a privacy policy & Chapter~A: absence as outcome with full denominator (\pa{sunyaev2015}, \pa{oloughlin2019}, \pa{tangari2021}) & Draft eligibility says ``reachable policy''; the meeting decided to include and report\\
D6 & Native and platform health apps & Chapter~S: scope by function, not vendor; no corpus paper treats platform apps separately & Needs a rule and a stated reason\\
D7 & Language rule per material type & Chapter~A and D: reviewer-competence rule (\pa{schultchen2021}); language match of policy and app (\pa{alfawzan2022}); original wording coded & English first, German fallback differs from the draft's German-first corpus\\
D8 & Two-coder procedure & Chapters~C and I: subset share, statistic, threshold, independence (\pa{oconnor2020}, \pa{lombard2002}, \pa{hallgren2012}, \pa{campbell2013}); single coder as limitation (\pa{mcdonald2019}) & Who the second coder is; the primary statistic and AC1 sensitivity\\
D9 & LLM role and acceptance thresholds & Chapter~L: human-first workflow (\pa{pangakis2023}); checklist (\pa{tornberg2024}); pre-registration (\pa{baumann2025}); not for inductive coding (\pa{mayring2025}) & Thresholds per code and for rare codes; model choice (open weights)\\
D10 & Walkthrough checklist and outputs & Chapters~W and C: CW questions (\pa{wharton1994}, \pa{spencer2000}); registration--use--leaving (\pa{light2018}); anchored levels (\pa{hutton2018}); click counts (\pa{habib2019}) & Adaptation named; outputs per step; one inspector or two\\
D11 & Device scope (apps only) & Chapter~S and D: behaviour and traffic analysis excluded (\pa{arkalakis2024}, \pa{tangari2021} as contrast); wearables defined in \pa{niksirat2024} & State consequences for claims (communication only)\\
D12 & Basis of the harm tree (``best app'') & Chapter~H: reference trees (\pa{knabenhans2026}), harm types (\pa{citron2021}), contextual severity (\pa{brooks2017}) & Supervisor to confirm what ``best'' means; counts of applicable pathways are not a harm ranking\\
\bottomrule
\end{longtable}

\section{Where the corpus agrees}
\begin{itemize}
\item Document the search and selection as a counted cascade with reasons and dates; do not tune thresholds after seeing results.
\item Fix tasks, units, codebook and statistics before the main coding; keep a change log and freeze a release.
\item Report agreement per code, with raw agreement and marginals beside a chance-corrected value.
\item Validate any model-assisted coding against human labels on untouched data and report every configuration tried.
\item Describe communication as communication: store labels and policies are claims, not evidence of practice.
\end{itemize}

\section{Where the corpus disagrees}
\begin{itemize}
\item Percent agreement and which chance-corrected coefficient to prefer (\pa{lombard2002} against \pa{krippendorff2004}; \pa{feinstein1990}, \pa{gwet2008}).
\item Whether reliability statistics fit interpretive coding at all (\pa{braun2021} against \pa{oconnor2020}, \pa{mcdonald2019}).
\item Whether popularity or recent review volume is the better ranking basis, and whether either relates to quality (\pa{sunyaev2015}, \pa{alfawzan2022}, \pa{hutton2018}, \pa{malki2024}).
\item Whether an LLM may substitute for a second human coder (\pa{gilardi2023}, \pa{pangakis2023}, \pa{baumann2025}, \pa{mayring2025}).
\end{itemize}

\section{Gaps in the corpus}
\begin{itemize}
\item No source validates a cognitive-walkthrough adaptation for privacy information in apps (\pa{mahatody2010} finds thin validation of CW variants generally).
\item No source tests whether an app-store ranking from a search API equals the visible ranking.
\item No German-language privacy-policy corpus with a published codebook; German evidence comes from \pa{schultchen2021}, \pa{simon2023}, \pa{messner2020} (app quality) and \pa{bfarm2023}, \pa{wp292018} (law and guidance).
\item Evidence for harms in the health-app domain is mostly from surveys and policy reviews (\pa{gabriele2020}, \pa{cao2024}, \pa{moniz2026}, \pa{niksirat2024}); there is no prevalence study of harm \emph{communication} apart from single items such as \pa{zych2026}.
\item Candidate literature to add (see \texttt{candidates.md}): APA framework papers, Zimmeck et al.\ APP-350, Cicchetti and Feinstein 1990 Part II, Krippendorff 2011/2016, Mehrnezhad and Almeida 2021.
\end{itemize}
